\chapter{结论与未来工作}\label{chap:conclusion}

\section{主要研究结论}
本文围绕 PostgreSQL Bug 修复演化给出系统研究：
\begin{enumerate}
    \item \textbf{修复模式一致性}：约 60\% 以上的逻辑修复遵循“增加防护分支”或“显式化错误路径”模式。
    \item \textbf{验证模型有效性}：证明了将形式化验证引入主流开源项目分析是高效的，特别是在处理算术溢出方面。
    \item \textbf{核心模块压力}：存储与复制模块仍是修复高地，其代码 Churn 值反映了长期的稳定性挑战。
\end{enumerate}

\section{工程实践建议}
建议在工程流程中落地三项实践：
\begin{itemize}
    \item 建立修复模式知识库，在代码审查阶段自动识别“脆弱结构”；
    \item 推行轻量级形式化检查，对高危逻辑要求附带 Z3 验证脚本；
    \item 增强回归测试的语义断言，捕捉静默失效。
\end{itemize}

\section{后续工作展望}
后续可从三方面扩展：基于历史模式训练 LLM 补丁生成模型；研发能够从 C 源码直接抽象出 Z3 谓词的自动化工具；以及将本方法扩展至 MySQL 等其他主流数据库进行对比研究。